%% toc: P1 3 3.7,3.8
\section{The Ultrafilter Lemma as a theorem at its first statement}
\label{sec:p01-w02b-p01-005-ultrafilter-lemma-entry}

\emph{Purpose.} Definition~\ref{def:ultrafilter-lemma} in
\cref{sec:set-prerequisites-06} introduces the Ultrafilter Lemma inside a
definition environment. That environment gives the principle no standing as a
result, and its sentence asserting that the principle is strictly weaker over
\(\mathsf{ZF}\) than the Axiom of Choice carries neither a reference nor the
consistency hypothesis under which such a claim can be a theorem
\cite[ch.~``Sets, Relations, Functions, and Cardinality'', sec.~``Choice
principles used later'', definition ``Ultrafilter Lemma'']{FrozenFormalLogicUltrafilterLemma}.
This section restates the same principle at the same point of the text as a
numbered theorem with an explicit standing, and states the strictness claim in
its correct relative form, resting on the theorem of Halpern and L\'evy
\cite{HalpernLevy1971}. The full development of the principle (vocabulary,
Boolean equivalents, side conditions for applying it, examples and limits)
follows in \cref{sec:p01-w02-p01-009-ultrafilter-lemma}; the theorem below is
one and the same result as parts~1, 2 and~4 of
Theorem~\ref{thm:p01-w02-p01-009-ultrafilter-lemma} there.

\emph{Inputs.} A set \(I\); a family \(\mathcal F\) of subsets of \(I\)
satisfying the four conditions displayed in the theorem; the Axiom of Choice,
\(\mathsf{AC}\), where a part names it; and, for the independence part, the
hypothesis that \(\mathsf{ZF}\) is consistent.
\emph{Conclusions.} A family \(\mathcal U\supseteq\mathcal F\) with the same
four conditions that contains exactly one of each subset of \(I\) and its
complement (part~1); the \(\mathsf{ZF}\)-provable implication from
\(\mathsf{AC}\) to that principle (part~2); and, if \(\mathsf{ZF}\) is
consistent, the \(\mathsf{ZF}\)-unprovability of the converse implication
(part~3).

\subsection{Assumptions}
\label{sec:p01-w02b-p01-005-assumptions}

The ambient theory is \(\mathsf{ZF}\), the classical set theory used as
metatheory in this chapter; \(\mathsf{AC}\) is assumed only where a statement
names it. The principle quantifies over all sets \(I\) and all families
\(\mathcal F\subseteq\mathcal P(I)\), so it is one sentence of the ambient set
theory, not a formula of a first-order theory under study.
\(\mathsf{BPI}\) is the Boolean Prime Ideal principle of
Definition~\ref{def:p01-bpi-principles}: every nontrivial Boolean algebra has
a prime ideal. For a theory \(T\), \(\operatorname{Con}(T)\) is the assertion
that no contradiction is derivable from the axioms of \(T\); it is used here
only as a hypothesis about provability, as in
Theorem~\ref{thm:p01-bpi-choice-separation}. The four conditions on
\(\mathcal F\) are labelled (F1)--(F4) and the condition that a family
contains exactly one of each subset and its complement is labelled (D), as in
\cref{sec:p01-w02-p01-009-assumptions}.

\subsection{The statement and its standing}
\label{sec:p01-w02b-p01-005-statement}

\begin{hbtheorem}{thm:p01-w02b-p01-005-ultrafilter-lemma}{proved=annex:p01-w02b-p01-005-ultrafilter-lemma}{Ultrafilter Lemma, with its standing}
\begin{enumerate}
  \item \emph{(The lemma.)} Assume \(\mathsf{ZF}+\mathsf{AC}\). Let \(I\) be
    a set and let \(\mathcal F\subseteq\mathcal P(I)\) satisfy
    \[
      I\in\mathcal F,\quad\varnothing\notin\mathcal F,\quad
      X,Y\in\mathcal F\Rightarrow X\cap Y\in\mathcal F,
    \]
    and
    \[
      X\in\mathcal F\text{ and }X\subseteq Y\subseteq I
      \Rightarrow Y\in\mathcal F.
    \]
    Then there is \(\mathcal U\subseteq\mathcal P(I)\) containing
    \(\mathcal F\) that satisfies the same conditions and, for every
    \(X\subseteq I\), contains exactly one of \(X\) and \(I\setminus X\).
  \item \emph{(Choice suffices.)} \(\mathsf{ZF}\) proves the implication from
    \(\mathsf{AC}\) to the principle of part~1, stated for all \(I\) and all
    \(\mathcal F\).
  \item \emph{(Choice is not implied, relative to consistency.)} If
    \(\mathsf{ZF}\) is consistent, then \(\mathsf{ZF}\) does not prove the
    implication from the principle of part~1 to \(\mathsf{AC}\). This part is
    derived from the theorem of Halpern and L\'evy
    \[
      \operatorname{Con}(\mathsf{ZF})\ \Longrightarrow\
      \operatorname{Con}(\mathsf{ZF}+\mathsf{BPI}+\neg\mathsf{AC})
    \]
    \cite{HalpernLevy1971}, recorded as
    Theorem~\ref{thm:p01-bpi-choice-separation}, together with the
    \(\mathsf{ZF}\)-equivalence of the principle with \(\mathsf{BPI}\)
    (Theorem~\ref{thm:p01-w02-p01-009-ultrafilter-lemma}, part~3).
\end{enumerate}
Consequently, provided \(\mathsf{ZF}\) is consistent, the Ultrafilter Lemma is
strictly weaker over \(\mathsf{ZF}\) than the Axiom of Choice, and it is not a
second name for that axiom. It is a metatheoretic choice principle about
families of subsets, not a formula of a first-order theory under study.
\end{hbtheorem}

Parts~1, 2 and~3 are, in hypotheses, notation and conclusion, parts~1, 2
and~4 of Theorem~\ref{thm:p01-w02-p01-009-ultrafilter-lemma}; the two labels
name one result, stated here where the principle first appears and developed
in \cref{sec:p01-w02-p01-009-ultrafilter-lemma}. The proof in
Annex~\ref{annex:p01-w02b-p01-005-ultrafilter-lemma} derives parts~1 and~2
from results proved in this handbook and derives part~3 from one external
premise, the Halpern--L\'evy relative-consistency theorem, whose model
construction is not reproduced in this handbook.

\emph{Correspondence with the inherited definition.} Each clause of
Definition~\ref{def:ultrafilter-lemma}, which stays on the reading route
unchanged, is carried as follows. The index set \(I\), the two displayed
condition blocks on \(\mathcal F\) and the conclusion about \(\mathcal U\) are
part~1, with the source's power-set symbol rendered as \(\mathcal P(I)\). The
sentence calling the principle metatheoretic and about families of subsets is
the last sentence of the theorem. The strictness sentence and its consequence
that the principle is not a second name for Choice are part~2 together with
part~3 and the sentence beginning ``Consequently''; the corrections made here
are the added consistency hypothesis and the added citation. The announcement
that a later chapter supplies short terms for the two displayed families is
met by Lemma~\ref{thm:p01-w02-p01-009-terms}, which identifies them with
proper set filters and ultrafilters in the sense of
Definition~\ref{def:p01-bpi-filter-ideal}.

\subsection{Worked example}
\label{sec:p01-w02b-p01-005-example}

\begin{hbexample}{def:p01-w02b-p01-005-worked}{An infinite index set where the lemma leaves a choice open}
Let \(I=\mathbb N\), let \(E\) be the set of even numbers, and for
\(n\in\mathbb N\) write \([0,n)=\{0,\ldots,n-1\}\). Let
\(\mathcal F=\{\,Y\subseteq\mathbb N:E\setminus[0,n)\subseteq Y\text{ for
some }n\in\mathbb N\,\}\).

\emph{The hypotheses hold without Choice.} \(\mathbb N\in\mathcal F\) with
\(n=0\). Each \(E\setminus[0,n)\) is infinite, hence nonempty, so
\(\varnothing\notin\mathcal F\). If \(E\setminus[0,n)\subseteq Y\) and
\(E\setminus[0,m)\subseteq Z\), then with \(k\) the larger of \(n\) and
\(m\), \(E\setminus[0,k)\subseteq Y\cap Z\). Supersets inside
\(\mathbb N\) of a member of \(\mathcal F\) are in \(\mathcal F\) with the
same \(n\).

\emph{What every extension contains.} Let \(\mathcal U\) be any family given
by part~1 of Theorem~\ref{thm:p01-w02b-p01-005-ultrafilter-lemma}. It
contains \(E\) (take \(n=0\)). It contains no finite set \(X\): such an
\(X\) lies inside some \([0,n)\), so \(\mathbb N\setminus X\supseteq
E\setminus[0,n)\) belongs to \(\mathcal F\subseteq\mathcal U\), and (D)
excludes \(X\).

\emph{What the lemma leaves open.} Let \(M\) be the set of multiples of
\(4\) and \(M'=E\setminus M\). The family \(\mathcal F_M\) of sets containing
some \(M\setminus[0,n)\) satisfies (F1)--(F4) by the argument above with
\(M\) in place of \(E\), contains \(M\), and contains \(\mathcal F\) because
\(M\setminus[0,n)\subseteq E\setminus[0,n)\). Applying part~1 to
\(\mathcal F_M\) gives a family \(\mathcal U\supseteq\mathcal F\) with
\(M\in\mathcal U\). The same steps with \(M'\) give
\(\mathcal U'\supseteq\mathcal F\) with \(M'\in\mathcal U'\); since
\(M'\subseteq\mathbb N\setminus M\), condition (F4) puts
\(\mathbb N\setminus M\) in \(\mathcal U'\), and (D) keeps \(M\) out of it.
So the four conditions on \(\mathcal F\) do not determine whether \(M\)
belongs to a deciding extension; the lemma asserts existence only, and each of
the two applications here used it, through \(\mathsf{AC}\), once. For a
finite index set no choice principle is needed
(Corollary~\ref{thm:p01-bpi-set-boundaries};
\cref{def:p01-w02-p01-009-worked}).
\end{hbexample}

\subsection{Limits of the strictness claim}
\label{sec:p01-w02b-p01-005-limits}

\begin{hbwarning}{def:p01-w02b-p01-005-strictness-reading}{Reading the inherited strictness sentence}
\begin{enumerate}
  \item Without its hypothesis the strictness sentence of
    Definition~\ref{def:ultrafilter-lemma} is not a theorem of anything
    weaker than the consistency of \(\mathsf{ZF}\): if \(\mathsf{ZF}\) were
    inconsistent, it would prove every sentence, including the implication
    from the lemma to \(\mathsf{AC}\). Part~3 of
    Theorem~\ref{thm:p01-w02b-p01-005-ultrafilter-lemma} is the form in which
    the claim holds.
  \item Part~3 is a claim about provability in \(\mathsf{ZF}\). It does not
    assert that Choice is false, and it does not assert that \(\mathsf{ZF}\)
    alone fails to prove the lemma; it compares the lemma with Choice only.
  \item The model of \(\mathsf{ZF}+\mathsf{BPI}+\neg\mathsf{AC}\) on which
    part~3 rests is cited from Halpern and L\'evy \cite{HalpernLevy1971} and
    is not constructed in this handbook. The derivation of part~3 from that
    cited theorem is complete in the annex.
  \item Part~1 is stated under \(\mathsf{AC}\), which suffices; by part~3, if
    \(\mathsf{ZF}\) is consistent it cannot prove the principle of part~1
    to imply \(\mathsf{AC}\). Over \(\mathsf{ZF}\) the principle of part~1
    is equivalent to \(\mathsf{BPI}\) and to propositional compactness for an
    arbitrary set of variables
    (Theorem~\ref{thm:p01-w02-p01-009-ultrafilter-lemma}, part~3), so in
    \(\mathsf{ZF}+\mathsf{BPI}\) the conclusion of part~1 holds without
    \(\mathsf{AC}\). The further limits are listed in
    \cref{def:p01-w02-p01-009-limits}.
\end{enumerate}
\end{hbwarning}

\subsection{Relationships and sources}
\label{sec:p01-w02b-p01-005-relationships}

The development that this theorem opens is
\cref{sec:p01-w02-p01-009-ultrafilter-lemma}: the names for the two families
(Lemma~\ref{thm:p01-w02-p01-009-terms}), the full theorem with the
equivalents (Theorem~\ref{thm:p01-w02-p01-009-ultrafilter-lemma}), the side
conditions for applying the lemma to a given family
(Lemma~\ref{thm:p01-w02-p01-009-side-conditions}) and the procedure built on
them (\cref{def:p01-w02-p01-009-procedure}). The Boolean development is
\cref{sec:p01-boolean-prime-ideal-and-ultrafilters}, including the Choice
implication (Theorem~\ref{thm:p01-bpi-ac-implication}), the four-way
equivalence (Theorem~\ref{thm:p01-bpi-equivalence}) and the cited separation
(Theorem~\ref{thm:p01-bpi-choice-separation}). The Zorn step is
Theorem~\ref{thm:p06-ac-to-zorn}.

The statement of part~1, its displayed conditions and its two qualifications
follow the inherited manuscript's definition
\cite[ch.~``Sets, Relations, Functions, and Cardinality'', sec.~``Choice
principles used later'', definition ``Ultrafilter Lemma'']{FrozenFormalLogicUltrafilterLemma}.
The relative-consistency theorem behind part~3 is that of Halpern and L\'evy
\cite{HalpernLevy1971}. The proofs in the annex and the worked example are
editorial arguments written for this handbook.
